\documentclass[12pt]{article}
\usepackage{amsmath,amssymb}
\usepackage{hyperref}
\usepackage{url}

\title{Status Report: \texttt{erdos\_problems-1} --- Mis-targeted Placeholder Issue
(no determinable statement), with corrected accessibility findings}
\author{solver-agent}
\date{}

\begin{document}
\maketitle

\section*{Summary}

Issue \texttt{erdos\_problems-1} carries \textbf{no determinable mathematical
statement}, so no honest proof can be written \emph{for this particular issue}.
However, the earlier conclusion that ``no Erd\H{o}s problem statement is
accessible through any channel'' was \textbf{incorrect}: the statements are in
fact fully accessible, and the genuine per-problem proof tasks exist as
\emph{separate} issues. The defect in \texttt{erdos\_problems-1} is that it is a
\emph{dataset-level placeholder}, not a pointer to any single problem.

\section*{Corrected accessibility findings (this run)}

The previous reviews queried the API at the bare path \texttt{/api/\dots},
which $404$s because the SPA is served behind a proxy prefix. The frontend's own
shim reveals the correct base path \texttt{/rmac/solve}. Using it:

\begin{itemize}
\item \texttt{GET /rmac/solve/api/ds/problems?dataset=erdos\_problems} returns a
paginated list of the individual problems, with ids of the form
\texttt{erdos\_0001}, \texttt{erdos\_0060}, \texttt{erdos\_0108}, \dots
\item \texttt{GET /rmac/solve/api/ds/problem/erdos\_problems/erdos\_0001} returns
the \emph{verbatim statement}, e.g. Erd\H{o}s Problem \#1 (distinct subset sums):
if $A\subseteq\{1,\dots,N\}$ with $\lvert A\rvert=n$ has all subset sums
distinct, then $N\gg 2^{n}$ (status: \textbf{open}, \$500 prize).
\item \texttt{GET /rmac/solve/api/issues/erdos\_0001?dataset=erdos\_problems}
returns a \emph{properly titled} proof-task issue \texttt{erdos\_0001-1}
(``Proof: Erd\H{o}s Problem \#1'') whose system comment embeds the full
statement. The same holds for every individual problem
(\texttt{erdos\_0060-1}, \texttt{erdos\_0108-1}, \texttt{erdos\_0996-1}, \dots).
\end{itemize}

So the verbatim statements and the legitimate per-problem proof issues do exist
and are reachable.

\section*{Why \texttt{erdos\_problems-1} itself is still unresolvable}

The assigned issue is keyed to \texttt{problem\_id = erdos\_problems}, i.e. the
\emph{dataset slug}, not to any \texttt{erdos\_NNNN}. Concretely:

\begin{itemize}
\item Its title is ``Proof: (untitled)'', area ``(unspecified)'', and its body
and comments contain no statement, no problem number, no quantifiers, no
constants, and no status. (Contrast the per-problem issues, each of which embeds
its statement.)
\item The issue record exposes the fields
\texttt{id, problem\_id, title, status, labels, created\_at, created\_by,
comments, documents}; \textbf{none} links it to a specific \texttt{erdos\_NNNN}.
\item The suffix ``\,-1\,'' is merely the per-issue index (the first issue of a
given key). It is \emph{not} a problem number: the genuine first problem's issue
is \texttt{erdos\_0001-1}, distinct from \texttt{erdos\_problems-1}.
\end{itemize}

Consequently, selecting any concrete problem (e.g.\ guessing ``Problem \#1''
from the suffix) would be an unfounded guess. The dataset has $1217$ candidates,
Erd\H{o}s statements are phrasing-sensitive, and the listed targets are
overwhelmingly \emph{open} research problems (e.g.\ \#1 carries a \$500 prize and
remains unsolved). Fabricating a claim and a ``proof'' would violate grounding
and faithfulness and would, with overwhelming probability, prove the wrong
thing or claim a false resolution of an open problem.

\section*{Conclusion and precise remediation}

\textbf{Verdict:} \texttt{erdos\_problems-1} is a mis-targeted slug-level
placeholder with no determinable claim; it should remain blocked
(\texttt{in\_progress}), not resolved. This is a \emph{specification/targeting}
defect, not a mathematical-difficulty gap, and (correcting the prior reviews)
\emph{not} an accessibility failure.

\textbf{Any one of the following unblocks a real proof effort:}
\begin{enumerate}
\item Re-point this issue's \texttt{problem\_id} to a specific
\texttt{erdos\_NNNN} (and set its title/body to that problem's verbatim
statement and status); or
\item Work the already-correct per-problem issues \texttt{erdos\_NNNN-1}
instead, each of which carries its statement; or
\item Add a \texttt{problem.tex} to this workspace with the exact target
statement and its intended status.
\end{enumerate}
Once a concrete target is pinned, the standard pipeline applies (formalize
$\to$ decompose quantifiers $\to$ choose strategy $\to$ full proof $\to$ gap
audit). Note, however, that the per-problem targets sampled here are
long-standing \emph{open} problems, so a complete proof would constitute a
research breakthrough rather than a routine exercise.

\end{document}
